% !TeX program = pdflatex
% =============================================================================
% Gruppenpuzzle: Anwendungen der Lorentzkraft
% UZH Praktikum I -- Kantonsschule im Lee, HS2026
% Klassen: 4fh + 4cdeg (SP4), gemeinsames Material -- Termin noch offen
% (spaetere Doppellektion nach Teil 1, siehe lorentzkraft-teil2-lektionsplan.tex).
%
% Inhaltliche Grundlage: gruppenpuzzle-lorentzkraft-teil2-plan.md (dieser
% Ordner) -- fuenf gleichwertige Themen (nicht sechs, siehe Begruendung dort):
% Gruppe A Massenspektrograph (Bainbridge, inkl. Wien-Filter als erste
% Stufe), Gruppe B Kreisbeschleuniger (Zyklotron + Synchrotron gemeinsam),
% Gruppe C MHD-Antrieb, Gruppe D Polarlichter, Gruppe E Hall-Effekt.
% Stammgruppe = ein Mitglied pro Thema (5 Personen), Expertengruppe = ein
% Mitglied pro Stammgruppe -- Anzahl Stammgruppen haengt von der
% Klassengroesse ab und ist bewusst nicht fest codiert (auf Wunsch
% verallgemeinert, urspruenglich fuer 15 SuS/3 Stammgruppen entworfen).
%
% Findet NACH dem (zweifachen) Fadenstrahlrohr-Experiment statt (Teil 1 und
% Einstieg von Teil 2). F=qvB, F=IlB, die Drei-Finger-Regel, r=mv/qB, die
% e/m_e-Bestimmung und die Schraubenbahn sind deshalb bereits bekannt und
% werden in den fuenf Themenboxen nur zitiert, nicht neu hergeleitet (siehe
% Plan-Datei, Abschnitt "Vorausgesetztes Vorwissen"). Neu ist in den
% Themenboxen jeweils nur das Kraeftegleichgewicht bei gekreuzten Feldern
% (Wien-Filter/Hall-Effekt), die Umlauffrequenz f=qB/(2 pi m) (Kreis-
% beschleuniger), die Antriebskraft F=BIL ueber den Strom (MHD-Antrieb) und
% der Spiegeleffekt/Verlustkegel (Polarlichter). Diese neuen Formeln werden
% jeweils direkt im Fliesstext der Themenbox hergeleitet (Skript-Stil:
% erklaerender Text mit eingebetteten Gleichungen, bewusst NICHT in einer
% separaten reinen Formel-Box -- auf Wunsch). Jede Themenbox hat ausserdem
% einen Link zur passenden LEIFIphysik-Simulation zum Ausprobieren (Quelle:
% Unterricht Lorentzkraft Teil 2.docx, dieser Ordner). Jede Kontrollfrage
% hat neben der Verstaendnisfrage auch eine Rechenaufgabe bzw.\ Herleitung.
%
% Boxen-Stil, Makros (theoriebox, taskbox, groupbox, answerbox, \Bild,
% \Kopfzeile) 1:1 uebernommen aus gruppenpuzzle3-u-i-r.tex
% (elektrizitaet/u-i-r/), inkl. Loesungen sichtbar (\printanswers) fuer die
% Selbstkontrolle der Expertengruppen.
%
% Bilder: siehe Bilder/-Ordner (wien-filter.pdf, mass-spectrograph.pdf,
% cyclotron-function-static.pdf, cyclotron_PSI.png,
% Particles-synchrotron-diagram-alternating-gradient-ring-energies-accelerator.png
% (aus Literatur/...jpg mit Inkscape konvertiert -- Original war ein als
% .jpg benanntes GIF, das pdflatex nicht direkt einbinden kann),
% mhd-ship-yamato-1.jpg, Nordlichter-Beitragsbild.jpg,
% charge-helix-trajectory-static.pdf,
% csm_20231218_Polarlicht_iStockWdP_6d45db0da8.png, hall-voltage.pdf,
% blut_messung.png). Auf Wunsch NICHT mehr verwendet: magnetosphere.jpg,
% magnetic-mirror-static.pdf, aurora-borealis.jpg (Polarlichter-Thema, drei
% der bisherigen fuenf Bilder ersetzt/entfernt), sowie mhd-accelerator.pdf
% und mhd_skizze.png (MHD-Antrieb: Skizzen entfernt, nur noch das
% Yamato-1-Foto -- die SuS zeichnen die Skizze stattdessen selbst in der
% Kontrollfrage). Linearbeschleuniger-Thema
% (Kontrastfrage, linac.png) auf Wunsch entfernt -- die Aufgabe bleibt nur
% eine Vertiefungsaufgabe (Massenspektrograph+Hall-Effekt) plus die
% Synthesetabelle.
% =============================================================================
\documentclass[addpoints,11pt,a4paper]{exam}

\usepackage[T1]{fontenc}
\usepackage[utf8]{inputenc}
\usepackage[ngerman]{babel}
\usepackage{lmodern}
\usepackage[a4paper,margin=2.2cm,headheight=14pt]{geometry}
\usepackage{amsmath}
\usepackage{siunitx}
% Hausstil: zusammengesetzte Einheiten immer als echter Bruch (\frac), nicht
% als Inline-Schraegstrich oder \tfrac -- gilt global fuer jedes \si/\SI.
\sisetup{per-mode=fraction,fraction-command=\frac}
\usepackage{array,booktabs,tabularx,multirow}
\usepackage{enumitem}
\usepackage{xcolor}
\usepackage{colortbl}
\usepackage{tikz}
\usepackage{physics}
\usepackage[most]{tcolorbox}
\usepackage{lastpage}
\usepackage{graphicx}
\usepackage{hyperref}

\usetikzlibrary{arrows.meta,calc,angles,quotes}

\usepackage{setspace}
\usepackage{needspace}
\setlength{\parindent}{0pt}
\setlength{\parskip}{0.6em}
\setstretch{1.08}
\setlist[itemize]{itemsep=0.4em}

% -----------------------------
% Farben (Corporate-Farbschema Physik KFR)
% -----------------------------
\definecolor{titleblue}{HTML}{183A6B}
\definecolor{accentblue}{HTML}{2E6F9E}
\definecolor{lightblue}{HTML}{EAF4FB}
\definecolor{lightgray}{HTML}{F4F6F8}
\definecolor{darkgray}{HTML}{4A4A4A}
\definecolor{accentorange}{HTML}{D9720C}
\definecolor{accentpurple}{HTML}{7B2D8E}
\definecolor{lightorange}{HTML}{FCEEE0}
\definecolor{lightpurple}{HTML}{F3EAF7}

\hypersetup{
  colorlinks=true,
  urlcolor=accentblue,
  linkcolor=accentblue,
  pdftitle={Gruppenpuzzle: Anwendungen der Lorentzkraft},
  pdfauthor={Loris Keller}
}

% -----------------------------
% Spaltentypen
% -----------------------------
\newcolumntype{Y}{>{\centering\arraybackslash}X}
\newcolumntype{P}[1]{>{\raggedright\arraybackslash}p{#1}}

% -----------------------------
% Boxen
% -----------------------------
\tcbset{before skip=10pt, after skip=10pt}
\newtcolorbox{titlebox}{
  enhanced,
  colback=titleblue,
  colframe=titleblue,
  boxrule=0pt,
  arc=3mm,
  left=3mm,right=3mm,top=3mm,bottom=3mm
}

% Praktikum I: Lernziele/Einleitung/Theorie/Aufgaben werden NICHT farbig
% gerahmt (nur Formel-, Hinweis- und Antwortbox bleiben Boxen) -- siehe
% institutionen/UZH-Praktikum-I/CLAUDE.md, Abschnitt "Boxen-Layout".
\newenvironment{infobox}{%
  \par\addvspace{10pt}\noindent
}{%
  \par\addvspace{10pt}
}

% Titel von Theorie-/Aufgaben-/Gruppenarbeit-Abschnitten: groesser, farbig
% und mit duenner Trennlinie statt schlichtem Fettdruck.
\newcommand{\Teiltitel}[2]{%
  \noindent{\Large\bfseries\color{#1}#2}\par
  \vspace{1pt}\noindent{\color{#1}\rule{\linewidth}{0.8pt}}\par\smallskip\noindent
}

\newenvironment{theoriebox}[1][Theorie]{%
  \par\addvspace{12pt}\Teiltitel{titleblue}{#1}
}{%
  \par\addvspace{10pt}
}

\newtcolorbox{taskbox}[1][]{
  enhanced, breakable,
  colback=white, colframe=white, boxrule=0pt,
  left=0mm,right=0mm,top=0mm,bottom=0mm,
  before skip=12pt, after skip=10pt,
  before upper={%
    \def\tmpboxtitle{#1}%
    \ifx\tmpboxtitle\empty\else\Teiltitel{accentblue}{#1}\fi
  }
}

\newtcolorbox{groupbox}[1][Gruppenarbeit]{
  enhanced, breakable,
  colback=white, colframe=white, boxrule=0pt,
  left=0mm,right=0mm,top=0mm,bottom=0mm,
  before skip=12pt, after skip=10pt,
  before upper={\Teiltitel{accentorange}{#1}}
}

\newenvironment{beispielbox}[1][Beispiel]{%
  \par\addvspace{10pt}\noindent\textbf{#1}\par\smallskip\noindent
}{%
  \par\addvspace{10pt}
}

\newtcolorbox{hinweisbox}[1][Hinweis]{
  enhanced,
  breakable,
  colback=lightorange,
  colframe=accentorange,
  title={#1},
  fonttitle=\bfseries,
  boxrule=0.7pt,
  arc=2mm,
  left=5mm,right=5mm,top=2mm,bottom=2mm
}

\newtcolorbox{formelbox}[1][Formel]{
  enhanced,
  breakable,
  colback=lightpurple,
  colframe=accentpurple,
  title={#1},
  fonttitle=\bfseries,
  boxrule=0.9pt,
  arc=2mm,
  left=5mm,right=5mm,top=2mm,bottom=2mm
}

\newtcolorbox{answerbox}{
  breakable,
  colback=white,
  colframe=gray!55,
  boxrule=0.5pt,
  arc=1.5mm,
  left=2mm,right=2mm,top=2mm,bottom=2mm
}

% Bild-Helfer: zentriertes Bild mit spuerbarem Abstand zum umgebenden Text.
% #1 = Breite, #2 = Dateipfad.
\newcommand{\Bild}[2]{%
  \par\vspace{0.6cm}
  \begin{center}
    \includegraphics[width=#1]{#2}
  \end{center}
  \vspace{0.6cm}
}

% --- Loesungen ein-/ausblenden -----------------------------------------------
%\noprintanswers
\printanswers

\pointformat{}
\renewcommand{\solutiontitle}{\noindent\color{titleblue}\textbf{L\"osung:}\enspace}

\pagestyle{headandfoot}
\cfoot{\textcolor{darkgray}{Seite \thepage\ von \pageref{LastPage}}}

\newcommand{\Kopfzeile}[4]{%
  \begin{titlebox}
    {\color{white}\sffamily\bfseries\Large #4}
  \end{titlebox}
  \vspace{0.5em}
  \noindent\textbf{Klasse:}~#1 \hfill \textbf{Datum:}~#2 \hfill \textbf{Thema:}~#3
    \hfill \textbf{Lehrperson:}~Loris Keller
  \vspace{0.8em}
  \lhead{\textcolor{darkgray}{#3}}
  \rhead{\textcolor{darkgray}{#4}}
}

\newlength{\antwortraumlen}
\newcommand{\Antwortraum}[1]{%
  \pgfmathsetlength{\antwortraumlen}{#1*2.5cm}%
  \begin{answerbox}
    \rule{0pt}{\antwortraumlen}%
  \end{answerbox}%
}

% Werte fuer Kopfzeile zentral setzen
\newcommand{\Klasse}{4fh / 4cdeg (SP4)}
\newcommand{\Datum}{Termin noch offen (nach Teil 2)}
\newcommand{\Thema}{Lorentzkraft (Teil 2): Gruppenpuzzle}
\newcommand{\ArbeitsblattTitel}{Gruppenpuzzle: Anwendungen der Lorentzkraft}

\begin{document}

\Kopfzeile{\Klasse}{\Datum}{\Thema}{\ArbeitsblattTitel}

% =============================================================================
% Einfuehrung: Was ist ein Gruppenpuzzle, Ablauf
% =============================================================================
\textbf{Gruppenpuzzle: Anwendungen der Lorentzkraft}

Dieses Arbeitsblatt bearbeiten Sie nicht allein, sondern als
\textbf{\emph{Gruppenpuzzle}} (auch Jigsaw genannt): Der Stoff wird in fünf
gleichwertige Themen aufgeteilt. Zuerst wird jede und jeder von Ihnen
Expertin oder Experte für genau eines dieser fünf Themen, danach bringen Sie
Ihr Wissen in eine neue Gruppe ein, in der alle fünf Themen zusammenkommen.

\textbf{Ablauf:}
{\sloppy
\begin{enumerate}[leftmargin=1.2cm,itemsep=0.3em]
  \item \textbf{\emph{Einteilung}} (ca.\ 3 Min.): Sie arbeiten in
  Fünfergruppen, den \textbf{\emph{Stammgruppen}}. Jedes Mitglied Ihrer
  Stammgruppe erhält eines der fünf Themen unten: Gruppe A
  (Massenspektrograph), Gruppe B (Kreisbeschleuniger), Gruppe C
  (MHD-Antrieb), Gruppe D (Polarlichter) oder Gruppe E (Hall-Effekt).
  \item \textbf{\emph{Erarbeitung (Expertenphase)}} (ca.\ 14 Min.): Alle mit
  demselben Buchstaben (derselben Gruppe A--E) setzen sich zusammen
  (\textbf{\emph{Expertengruppe}}) und lesen gemeinsam ihren Abschnitt.
  Lösen Sie darin auch die Kontrollfrage Ihrer Gruppe.
  \item \textbf{\emph{Vermittlung (Austauschphase)}} (ca.\ 14 Min.): Sie
  kehren in Ihre Stammgruppe zurück. Reihum erklärt jede und jeder ihr/sein
  Thema den anderen vier Gruppenmitgliedern, bis alle fünf Themen einmal
  erklärt wurden.
  \item \textbf{\emph{Vertiefungsphase}}: Am Schluss löst Ihre komplette
  Stammgruppe gemeinsam die drei Aufgaben ganz am Ende dieses
  Arbeitsblatts, die mehrere der fünf Themen miteinander verknüpfen.
\end{enumerate}
\par}

Alle fünf Themen bauen auf der Lorentzkraft $F=qvB$ auf, die Sie am
Fadenstrahlrohr bereits kennengelernt haben. \textbf{Bearbeiten Sie in der
Expertenphase nur den Abschnitt Ihrer eigenen Gruppe.}

% =============================================================================
% Gruppe A: Massenspektrograph (Bainbridge)
% =============================================================================
\needspace{14cm}
\begin{theoriebox}[Gruppe A: Massenspektrograph (Bainbridge)]
Der \textbf{\emph{Bainbridge-Massenspektrograph}} bestimmt die Masse $m$
geladener Ionen (bekannte Ladung $q$) in zwei Stufen.

\textbf{Stufe 1 -- Geschwindigkeitsfilter (Wien-Filter):} Die Ionen
durchqueren zunächst ein Gebiet mit gekreuzten $E$- und $B_1$-Feldern (beide
senkrecht zueinander und senkrecht zur Bewegungsrichtung $v$). Die bereits
bekannte Lorentzkraft $qvB_1$ wirkt hier einer elektrischen Kraft $qE$
\emph{gleichzeitig und entgegengesetzt} entgegen -- neu ist dieses
\textbf{\emph{Kräftegleichgewicht}} selbst. Nur Ionen, bei denen sich beide
Kräfte exakt aufheben, fliegen geradeaus durch die Lochblende; schnellere
oder langsamere Ionen werden seitlich abgelenkt und damit ausgefiltert. Im
Gleichgewicht gilt
\[
  F_{el}=F_L \quad\Longrightarrow\quad qE=qv_0B_1.
\]
Die Ladung $q$ steht auf beiden Seiten und kürzt sich heraus, sodass die
selektierte Geschwindigkeit nur noch von den beiden Feldern abhängt:
\[
  v_0=\frac{E}{B_1}.
\]
Alle den Filter verlassenden Ionen haben also exakt dieselbe, aus $E$ und
$B_1$ bekannte Geschwindigkeit $v_0$, unabhängig von ihrer Masse.

\Bild{7.5cm}{Bilder/wien-filter.pdf}

\textbf{Stufe 2 -- Massentrennung im Magnetfeld:} Diese
geschwindigkeitsselektierten Ionen treten anschliessend in ein zweites,
reines Magnetfeld $B_2$ ein, wo wieder die bereits bekannte Formel
$r=\frac{mv_0}{qB_2}$ gilt, wie am Fadenstrahlrohr bereits gezeigt. Löst
man diese Gleichung nach der gesuchten Masse auf, multipliziert man zuerst
beide Seiten mit $qB_2$ und teilt anschliessend durch $v_0$:
\[
  r=\frac{mv_0}{qB_2}
  \quad\Longrightarrow\quad
  rqB_2=mv_0
  \quad\Longrightarrow\quad
  m=\frac{qB_2r}{v_0}.
\]
Schwerere Ionen laufen auf grösserem Radius und treffen weiter vom
Eintrittspunkt entfernt auf den Schirm auf; aus dem gemessenen Abstand
lässt sich bei bekanntem $v_0$, $B_2$, $q$ die Masse $m$ bestimmen --
historisch die Methode, mit der \textbf{\emph{Isotope}} (Atome desselben
Elements mit unterschiedlicher Masse) entdeckt und getrennt wurden. Ohne
die Geschwindigkeitsfilter-Stufe wäre die Massentrennung ungenau, weil $r$
sowohl von $m$ als auch von $v$ abhängt -- erst ein für alle Ionen
identisches $v_0$ macht $r$ zu einem eindeutigen Mass für $m$.

\Bild{7.5cm}{Bilder/mass-spectrograph.pdf}

\textbf{Zum Ausprobieren:} die
\href{https://www.leifiphysik.de/elektrizitaetslehre/bewegte-ladungen-feldern/downloads/wienscher-geschwindigkeitsfilter-grundwissen-simulation}{Simulation
Wienscher Geschwindigkeitsfilter} für Stufe 1 im Detail, ergänzend der
interaktive
\href{https://www.leifiphysik.de/elektrizitaetslehre/bewegte-ladungen-feldern/versuche/bainbridge-massenspektrometer}{Versuch
Bainbridge-Massenspektrometer} (zeigt den vollständigen Aufbau, beide
Stufen).
\end{theoriebox}

\begin{questions}
\begin{taskbox}[Kontrollfrage: Massenspektrograph (Gruppe A)]
\question[5]
\begin{parts}
  \part[3] Wozu braucht der Bainbridge-Massenspektrograph den
  Geschwindigkeitsfilter (Stufe 1) überhaupt, wenn für die eigentliche
  Massenbestimmung in Stufe 2 ohnehin nur $r=mv/qB_2$ verwendet wird? Was
  ginge schief, wenn die Ionen mit unterschiedlichen Geschwindigkeiten in
  Stufe 2 einträten?
  \begin{answerbox}
  \vspace{2.6cm}
  \end{answerbox}
  \begin{solution}
  Ohne den Geschwindigkeitsfilter hätten die Ionen beim Eintritt in Stufe 2
  unterschiedliche Geschwindigkeiten $v$. Da $r=mv/qB_2$ sowohl von $m$ als
  auch von $v$ abhängt, könnte ein grosser gemessener Radius entweder von
  einer grossen Masse oder von einer hohen Geschwindigkeit stammen -- der
  Radius wäre kein eindeutiges Mass für $m$. Erst wenn alle Ionen dank des
  Filters dieselbe Geschwindigkeit $v_0$ haben, hängt $r$ nur noch von $m$
  ab, und die Masse lässt sich eindeutig bestimmen.
  \end{solution}

  \part[2] Ein Wien-Filter hat $E=\SI{8.0e4}{\volt\per\meter}$ und
  $B_1=\SI{0.20}{\tesla}$. Welche Geschwindigkeit $v_0$ wird selektiert?
  \begin{answerbox}
  \vspace{1.8cm}
  \end{answerbox}
  \begin{solution}
  $v_0=\dfrac{E}{B_1}=\dfrac{\SI{8.0e4}{\volt\per\meter}}{\SI{0.20}{\tesla}}
  =\SI{4.0e5}{\meter\per\second}$.
  \end{solution}
\end{parts}
\end{taskbox}
\end{questions}

% =============================================================================
% Gruppe B: Kreisbeschleuniger (Zyklotron und Synchrotron)
% =============================================================================
\needspace{14cm}
\begin{theoriebox}[Gruppe B: Kreisbeschleuniger -- Zyklotron und Synchrotron]
Beim \textbf{\emph{Zyklotron}} befindet sich das Teilchen in einem homogenen
Magnetfeld zwischen zwei D-förmigen Hohlelektroden (\glqq Dees\grqq{}) mit
einer Wechselspannung dazwischen. Bei jedem Durchqueren des Spalts wird es
durch das $E$-Feld beschleunigt, innerhalb der Dees läuft es feldfrei im
Kreis -- hier gilt wieder die bereits bekannte Formel $r=mv/qB$, wie am
Fadenstrahlrohr bereits gezeigt, nur wächst jetzt $v$ mit jeder
Beschleunigung, also wächst auch $r$ mit: Es entsteht eine Spiralbahn mit
wachsendem Radius statt eines einzelnen Kreises.

\Bild{8cm}{Bilder/cyclotron-function-static.pdf}

Neu ist die \textbf{\emph{Umlauffrequenz}} $f$, mit der das Teilchen die
Dees durchläuft. Die Umlaufzeit $T$ ist die Bahnlänge (der Kreisumfang
$2\pi r$) geteilt durch die Geschwindigkeit:
\[
  T=\frac{2\pi r}{v}.
\]
Setzt man hier die bereits bekannte Formel $r=mv/qB$ ein, kürzt sich die
Geschwindigkeit $v$ heraus:
\[
  T=\frac{2\pi}{v}\cdot\frac{mv}{qB}=\frac{2\pi m}{qB}.
\]
Die Frequenz ist der Kehrwert der Umlaufzeit, $f=1/T$, also
\[
  f=\frac{qB}{2\pi m}.
\]
Weil $v$ beim Einsetzen verschwunden ist, ist $f$ (nicht-relativistisch)
geschwindigkeitsunabhängig, weshalb eine konstante Wechselspannungs-Frequenz
genügt, obwohl das Teilchen laufend schneller wird. Bei sehr hohen Energien
versagt dieses Prinzip trotzdem (relativistische Massenzunahme, immer
grösserer Radius). Das
\textbf{\emph{Synchrotron}} löst das, indem $B$ synchron mit der wachsenden
Energie erhöht wird, sodass der Radius konstant bleibt: Die Teilchen laufen
auf einem festen Ring, werden an mehreren Beschleunigungskavitäten
nachbeschleunigt und von Ablenkmagneten auf der Kreisbahn gehalten.

\Bild{8.5cm}{Bilder/Particles-synchrotron-diagram-alternating-gradient-ring-energies-accelerator.png}

Bekanntestes Beispiel: der LHC am CERN. Ein reales, nahes Beispiel ist das
Zyklotron am Paul-Scherrer-Institut (PSI) in Villigen.

\Bild{7cm}{Bilder/cyclotron_PSI.png}

\textbf{Simulation zum Ausprobieren:}
\href{https://www.leifiphysik.de/elektrizitaetslehre/bewegte-ladungen-feldern/ausblick/zyklotron}{Zyklotron}
und
\href{https://www.leifiphysik.de/elektrizitaetslehre/bewegte-ladungen-feldern/ausblick/synchro-zyklotron-und-synchrotrone}{Synchro-Zyklotron
und Synchrotrone}.
\end{theoriebox}

\begin{questions}
\setcounter{question}{1}
\begin{taskbox}[Kontrollfrage: Kreisbeschleuniger (Gruppe B)]
\question[5]
\begin{parts}
  \part[3] Warum lassen sich mit einem Zyklotron nicht beliebig hohe
  Energien erreichen, mit einem Synchrotron aber schon? Was ist der
  entscheidende technische Unterschied zwischen den beiden?
  \begin{answerbox}
  \vspace{2.6cm}
  \end{answerbox}
  \begin{solution}
  Beim Zyklotron ist $B$ konstant. Wächst $v$ (bei sehr hohen Energien durch
  relativistische Massenzunahme auch $m$), wächst der Radius $r=mv/qB$
  laufend, und die Umlauffrequenz $f=qB/2\pi m$ passt irgendwann nicht mehr
  zur konstanten Wechselspannungs-Frequenz -- das Prinzip versagt. Das
  Synchrotron erhöht $B$ synchron mit der wachsenden Energie, sodass der
  Radius konstant bleibt: Die Teilchen laufen auf einem festen Ring, dadurch
  sind beliebig hohe Energien erreichbar (begrenzt nur durch Ringgrösse und
  Technik, nicht mehr durch das Prinzip selbst).
  \end{solution}

  \part[2] Berechnen Sie die Umlauffrequenz $f$ eines Protons
  ($q=\SI{1.6e-19}{\coulomb}$, $m=\SI{1.67e-27}{\kilogram}$) in einem
  Zyklotron mit $B=\SI{1.5}{\tesla}$.
  \begin{answerbox}
  \vspace{2cm}
  \end{answerbox}
  \begin{solution}
  $f=\dfrac{qB}{2\pi m}=\dfrac{\SI{1.6e-19}{\coulomb}\cdot\SI{1.5}{\tesla}}
  {2\pi\cdot\SI{1.67e-27}{\kilogram}}\approx\SI{2.3e7}{\hertz}=\SI{23}{\mega\hertz}$.
  \end{solution}
\end{parts}
\end{taskbox}
\end{questions}

% =============================================================================
% Gruppe C: MHD-Antrieb
% =============================================================================
\needspace{12cm}
\begin{theoriebox}[Gruppe C: MHD-Antrieb]
In einem mit Meerwasser (leitfähig, enthält Ionen) gefüllten Kanal wird über
Elektroden ein Strom $I$ quer zur Strömungsrichtung durch das Wasser
geschickt, während gleichzeitig ein Magnetfeld $B$ senkrecht dazu steht.
Auf die bewegten Ionen im Wasser wirkt dabei die bereits bekannte
Lorentzkraft $F_1=qv_dB$ ($v_d$: Driftgeschwindigkeit der Ionen). Um die
gesamte Antriebskraft zu erhalten, betrachten wir einen Kanalabschnitt der
Länge $L$ mit Querschnittsfläche $A$, in dem sich Ionen der Dichte $n$
(Anzahl pro Volumen) befinden. Darin befinden sich insgesamt
\[
  N=n\cdot A\cdot L
\]
Ionen, und jedes erfährt die Kraft $F_1$. Die Gesamtkraft ist die Summe
aller Einzelkräfte:
\[
  F=N\cdot F_1=(nAL)\cdot(qv_dB)=(nqv_dA)\cdot(LB).
\]
Der Ausdruck $nqv_dA$ ist gerade die über die Ladungsträgerdichte
definierte Stromstärke $I$ (Ladung pro Zeit durch den Querschnitt), sodass
\[
  F=BIL
\]
folgt -- dieselbe Formel, die Sie bereits als Kraft auf einen
stromdurchflossenen Leiter im Magnetfeld kennen, hier aber direkt aus der
Lorentzkraft auf die einzelnen Ionen hergeleitet. Diese Kraft wirkt als
\textbf{\emph{direkte Antriebskraft}} in eine dritte, zu $I$ und $B$
senkrechte Richtung und beschleunigt das Wasser aus dem Kanal -- das
Schiff wird angetrieben, ganz ohne bewegliche Teile (kein Propeller).

Reales Beispiel: die \textbf{\emph{Yamato-1}} (Japan, 1992), das erste
Schiff mit funktionierendem MHD-Antrieb.

\Bild{7cm}{Bilder/mhd-ship-yamato-1.jpg}

Trotz des funktionierenden Prinzips wird der MHD-Antrieb heute kaum
eingesetzt: Meerwasser leitet Strom nur schlecht, weshalb für eine
nennenswerte Kraft $F=BIL$ entweder ein sehr grosser Strom $I$ oder ein
sehr starkes Magnetfeld $B$ nötig ist. Ein grosser Strom durch das
schlecht leitende Wasser erzeugt dabei vor allem Wärme statt Antrieb, und
ein starkes Magnetfeld braucht schwere, teure supraleitende Magnete samt
Kühlung. Wie klein der Wirkungsgrad dadurch tatsächlich ausfällt, finden
Sie in der Kontrollfrage selbst heraus.

\textbf{Mehr dazu:}
\href{https://www.leifiphysik.de/elektrizitaetslehre/bewegte-ladungen-feldern/ausblick/mhd-generator-und-mhd-pumpe}{MHD-Generator
und MHD-Pumpe} (LEIFIphysik).
\end{theoriebox}

\begin{questions}
\setcounter{question}{2}
\begin{taskbox}[Kontrollfrage: MHD-Antrieb (Gruppe C)]
\question[12]
\begin{parts}
  \part[3] Zeichnen Sie eine beschriftete Skizze des MHD-Kanals von oben.
  Markieren Sie darin die Strömungsrichtung des Meerwassers (= Richtung
  der Kraft $F$), die Elektroden, die Richtung des Magnetfelds $B$ und die
  Richtung des Stroms $I$.
  \begin{answerbox}
  \vspace{5.5cm}
  \end{answerbox}
  \begin{solution}
  \begin{center}
  \begin{tikzpicture}[scale=1,>=Latex]
    \draw[thick] (0,0) rectangle (4,2);
    \foreach \x in {0.6,1.6,2.6,3.6}
      \foreach \y in {0.5,1.5}
        \node[gray] at (\x,\y) {$\times$};
    \node[gray] at (2,2.5) {$B$ (ins Blatt hinein)};
    \draw[very thick,red] (0,2) -- (4,2) node[midway,above]{Elektrode $+$};
    \draw[very thick,red] (0,0) -- (4,0) node[midway,below]{Elektrode $-$};
    \foreach \x in {1,2,3}
      \draw[-{Latex},thick,red] (\x,1.7) -- (\x,0.3);
    \node[red] at (3.85,1) {$I$};
    \draw[-{Latex},ultra thick,blue] (4.2,1) -- (5.8,1);
    \node[blue] at (5.8,1.4) {$F$ (Wasserfluss)};
  \end{tikzpicture}
  \end{center}
  $I$ fliesst von der Plus- zur Minuselektrode durchs Wasser, $B$ steht
  senkrecht dazu ins Blatt hinein; nach der Drei-Finger-Regel zeigt die
  Kraft $F$ (und damit die Wasserströmung) dann seitlich aus dem Kanal.
  \end{solution}

  \part[2] Woher kommt die Kraft, die das Meerwasser im MHD-Kanal
  beschleunigt, wenn es doch gar keine Schraube gibt? Nennen Sie die
  beiden Grössen, die dafür senkrecht zueinander stehen müssen, damit die
  Kraft in die gewünschte dritte Richtung wirkt.
  \begin{answerbox}
  \vspace{2.2cm}
  \end{answerbox}
  \begin{solution}
  Die Kraft ist die Lorentzkraft auf die im Meerwasser bewegten
  Ladungsträger, hier über den Strom ausgedrückt: $F=BIL$. Damit sie in
  die gewünschte dritte Richtung wirkt, müssen der Strom $I$ (quer durch
  das Wasser) und das Magnetfeld $B$ senkrecht zueinander stehen.
  \end{solution}

  \part[2] Durch einen MHD-Kanal mit Elektrodenabstand $L=\SI{1.0}{\meter}$
  fliesst ein Strom $I=\SI{1000}{\ampere}$ senkrecht zu einem Magnetfeld
  $B=\SI{0.50}{\tesla}$. Berechnen Sie die Antriebskraft $F$. Diese Formel
  kennen Sie bereits aus einem anderen Zusammenhang -- aus welchem?
  \begin{answerbox}
  \vspace{2.2cm}
  \end{answerbox}
  \begin{solution}
  $F=BIL=\SI{0.50}{\tesla}\cdot\SI{1000}{\ampere}\cdot\SI{1.0}{\meter}
  =\SI{500}{\newton}$. Die Formel kennen Sie bereits als Kraft auf einen
  stromdurchflossenen Leiter im Magnetfeld (Teil 1, z.\,B. beim
  Leiterschaukel-Experiment).
  \end{solution}

  \part[3] Bei $v=\SI{4.0}{\meter\per\second}$ Schiffsgeschwindigkeit
  leistet die Antriebskraft $F$ aus Teilaufgabe b) die mechanische
  Leistung $P_{mech}=Fv$. Die Elektroden werden dazu mit
  $U=\SI{25}{\volt}$ betrieben, sodass die aufgewendete elektrische
  Leistung $P_{el}=UI$ beträgt. Berechnen Sie $P_{mech}$, $P_{el}$ und
  daraus den Wirkungsgrad $\eta=P_{mech}/P_{el}$ (in Prozent).
  \begin{answerbox}
  \vspace{2.8cm}
  \end{answerbox}
  \begin{solution}
  $P_{mech}=Fv=\SI{500}{\newton}\cdot\SI{4.0}{\meter\per\second}=\SI{2000}{\watt}$.\\
  $P_{el}=UI=\SI{25}{\volt}\cdot\SI{1000}{\ampere}=\SI{25000}{\watt}$.\\
  $\eta=\dfrac{P_{mech}}{P_{el}}=\dfrac{\SI{2000}{\watt}}{\SI{25000}{\watt}}=0{,}08=\SI{8}{\percent}$.
  \end{solution}

  \part[2] Ein gewöhnlicher Propellerantrieb erreicht einen Wirkungsgrad
  von rund $\SI{70}{\percent}$. Vergleichen Sie mit Ihrem berechneten Wert
  aus Teilaufgabe c): Was bedeutet das für die Wirtschaftlichkeit des
  MHD-Antriebs im Vergleich zu einem herkömmlichen Schiff?
  \begin{answerbox}
  \vspace{2.2cm}
  \end{answerbox}
  \begin{solution}
  Der berechnete Wirkungsgrad von rund $\SI{8}{\percent}$ liegt weit unter
  dem eines Propellerantriebs ($\SI{70}{\percent}$). Für dieselbe
  Antriebsleistung braucht der MHD-Antrieb also fast neunmal mehr
  elektrische Energie -- er ist damit trotz des funktionierenden Prinzips
  wirtschaftlich klar unterlegen, was erklärt, wieso er sich gegenüber dem
  klassischen Propeller nie durchgesetzt hat.
  \end{solution}
\end{parts}
\end{taskbox}
\end{questions}

% =============================================================================
% Gruppe D: Polarlichter
% =============================================================================
\needspace{16cm}
\begin{theoriebox}[Gruppe D: Polarlichter]
Sie kennen dieses Beispiel schon aus Teil 1 -- jetzt schauen wir uns an,
wieso das genau so passiert. Der Sonnenwind (geladene Teilchen von der
Sonne) trifft auf die Magnetosphäre der Erde.

\Bild{7cm}{Bilder/Nordlichter-Beitragsbild.jpg}

\textbf{Wie die Spiralbahn entsteht:} Ein Teilchen, das schräg zu einer
Feldlinie einfliegt, hat eine Geschwindigkeitskomponente $v_\parallel$
entlang der Feldlinie und eine Komponente $v_\perp$ quer dazu -- die
bereits bekannte Schraubenbahn-Zerlegung. Nur $v_\perp$ erzeugt über die
Lorentzkraft eine Kreisbewegung um die Feldlinie; $v_\parallel$ bleibt
davon unbeeinflusst und trägt das Teilchen gleichzeitig die Feldlinie
entlang. Kreisbewegung quer plus gleichförmige Bewegung längs ergibt
zusammen die Spiralbahn, die die Feldlinie wie eine Feder umwickelt -- die
Feldlinie selbst ist dabei die \glqq Achse\grqq{} der Spirale.

\Bild{7cm}{Bilder/charge-helix-trajectory-static.pdf}
\Bild{6cm}{Bilder/polarlicht_illustration_flugbahn.png}

\textbf{Wie daraus die magnetische Flasche wird:} Neu ist, was passiert,
wenn das Teilchen auf diesem Spiralweg in ein stärker werdendes Feld
hineinläuft -- genau das passiert an den Polen, wo die Feldlinien der Erde
zusammenlaufen und $B$ deshalb zunimmt. Wie schon beim Fadenstrahlrohr
gilt: Die Lorentzkraft steht immer senkrecht zur Geschwindigkeit und kann
deshalb nur deren \emph{Richtung} ändern, nie ihren \emph{Betrag} -- die
Gesamtgeschwindigkeit $v=\sqrt{v_\parallel^2+v_\perp^2}$ des Teilchens
bleibt während der ganzen Bewegung konstant. Wird $B$ stärker, wächst also
$v_\perp$ genau so weit, wie $v_\parallel$ abnimmt, sodass die
Gesamtgeschwindigkeit gleich bleibt. Irgendwann ist die gesamte
Geschwindigkeit in $v_\perp$ umgewandelt, $v_\parallel=0$: An diesem
Umkehrpunkt wirkt die Lorentzkraft wie ein \textbf{\emph{Spiegel}} und
schickt das Teilchen die Feldlinie wieder zurück, Richtung Äquator und
weiter zum gegenüberliegenden Pol, wo sich derselbe Effekt wiederholt.

Ein Teilchen kann so wiederholt zwischen Nord- und Südpol hin- und
herpendeln, ohne je die Atmosphäre zu erreichen -- es ist in einer
natürlichen \textbf{\emph{magnetischen Flasche}} gefangen (den
Van-Allen-Gürteln), mit den beiden polnahen Zonen starken Feldes als
\glqq Flaschenhälse\grqq{}. Nur Teilchen, deren Geschwindigkeitsrichtung
beim Eintritt zu nahe an $v_\parallel$ liegt (zu wenig $v_\perp$,
sogenannter \textbf{\emph{Verlustkegel}}), erreichen den Umkehrpunkt gar
nicht, sondern schon vorher die dichte Atmosphäre -- dort erzeugen sie
durch Stösse mit Luftmolekülen das Polarlicht, bevorzugt in einem Ring
rund um die Pole (\textbf{\emph{Auroraloval}}), nicht direkt über den
Polen selbst, weil genau dort die Feldlinien am steilsten und am
dichtesten zusammenlaufen.

\Bild{6cm}{Bilder/csm_20231218_Polarlicht_iStockWdP_6d45db0da8.png}

\textbf{Mehr dazu:}
\href{https://www.leifiphysik.de/elektrizitaetslehre/bewegte-ladungen-feldern/ausblick/polarlicht}{Polarlicht}
und
\href{https://www.leifiphysik.de/elektrizitaetslehre/bewegte-ladungen-feldern/ausblick/teilchenbahnen-magnetfeldern}{Teilchenbahnen
in Magnetfeldern} (beide LEIFIphysik). \textbf{Aufgabe zum Vertiefen:}
\href{https://www.leifiphysik.de/elektrizitaetslehre/bewegte-ladungen-feldern/aufgabe/magnetische-flasche}{Magnetische
Flasche}.
\end{theoriebox}

\begin{questions}
\setcounter{question}{3}
\begin{taskbox}[Kontrollfrage: Polarlichter (Gruppe D)]
\question[8]
\begin{parts}
  \part[3] Ein Teilchen fliegt spiralförmig entlang einer Feldlinie in
  Richtung Nordpol, wo das Feld stärker wird. Was passiert dabei mit
  $v_\perp$ und $v_\parallel$, und wieso kehrt sich seine Bewegung
  irgendwann um, statt einfach weiterzufliegen?
  \begin{answerbox}
  \vspace{2.4cm}
  \end{answerbox}
  \begin{solution}
  Weil die Lorentzkraft nie Arbeit verrichtet, bleibt die gesamte
  Geschwindigkeit (Bewegungsenergie) des Teilchens konstant. Mit
  zunehmendem $B$ wächst $v_\perp$ auf Kosten von $v_\parallel$. Am
  Umkehrpunkt ist die gesamte Bewegungsenergie in $v_\perp$ umgewandelt
  ($v_\parallel=0$): Die Lorentzkraft wirkt dort wie ein Spiegel und
  schickt das Teilchen die Feldlinie wieder zurück.
  \end{solution}

  \part[2] Wieso entstehen Polarlichter trotzdem nicht überall auf der
  Erde gleich häufig, sondern vor allem in einem Ring rund um die Pole?
  \begin{answerbox}
  \vspace{1.8cm}
  \end{answerbox}
  \begin{solution}
  Nur Teilchen mit einer Geschwindigkeitsrichtung zu nahe an
  $v_\parallel$ (zu wenig $v_\perp$, Verlustkegel) erreichen den
  Umkehrpunkt gar nicht, sondern schon vorher die dichte Atmosphäre. Das
  geschieht bevorzugt in einem Ring rund um die Pole (Auroraloval), nicht
  direkt über den Polen selbst, weil dort die Feldlinien am steilsten und
  dichtesten zusammenlaufen.
  \end{solution}

  \part[3] Begründen Sie, wieso die Gesamtgeschwindigkeit
  $v=\sqrt{v_\parallel^2+v_\perp^2}$ eines Teilchens während der gesamten
  Bewegung entlang einer Feldlinie konstant bleibt, obwohl sich
  $v_\parallel$ und $v_\perp$ einzeln verändern. Nutzen Sie dabei, was Sie
  bereits vom Fadenstrahlrohr über die Richtung der Lorentzkraft wissen.
  \begin{answerbox}
  \vspace{2.4cm}
  \end{answerbox}
  \begin{solution}
  Die Lorentzkraft steht immer senkrecht zur Geschwindigkeit -- genau wie
  am Fadenstrahlrohr kann sie deshalb nur die Richtung der Bewegung
  ändern, nie deren Betrag. Die Gesamtgeschwindigkeit $v$ bleibt also
  konstant. Da $v^2=v_\parallel^2+v_\perp^2$ und $v$ konstant ist, muss
  $v_\perp$ zunehmen, wenn $v_\parallel$ abnimmt (und umgekehrt) -- die
  Summe der Quadrate bleibt gleich, auch wenn sich die einzelnen
  Komponenten gegenläufig verändern.
  \end{solution}
\end{parts}
\end{taskbox}
\end{questions}

% =============================================================================
% Gruppe E: Hall-Effekt
% =============================================================================
\needspace{14cm}
\begin{theoriebox}[Gruppe E: Hall-Effekt]
Bewegen sich Ladungsträger (z.\,B. in einem stromdurchflossenen Leiter oder
in einer strömenden leitfähigen Flüssigkeit wie Blut) quer zu einem
angelegten Magnetfeld $B$, wirkt auf sie die bereits bekannte Lorentzkraft
$qvB$ senkrecht zu $v$ und $B$ -- sie werden dadurch zu einer Seite hin
abgelenkt. Weil sie den Leiter/das Gefäss aber nicht verlassen können,
sammelt sich dort ein Ladungsüberschuss an, bis das dadurch entstehende
elektrische Feld $E_H$ die Lorentzkraft gerade kompensiert; neu ist dieses
\textbf{\emph{Kräftegleichgewicht}} selbst -- dieselbe Idee wie beim
Wien-Filter (Thema Massenspektrograph), nur mit vertauschter
Fragestellung: Dort ist $v_0$ vorgegeben und wird über $E$ und $B$
eingestellt; hier ist $v$ unbekannt und wird über die resultierende
Spannung gemessen. Im Gleichgewicht gilt
\[
  F_{el}=F_L \quad\Longrightarrow\quad qE_H=qvB
  \quad\Longrightarrow\quad E_H=vB,
\]
wobei die Ladung $q$ beim Teilen wieder herausgekürzt wurde.

\Bild{8cm}{Bilder/hall-voltage.pdf}

Die Feldstärke hängt über den Elektrodenabstand $d$ mit der messbaren
Spannung zusammen, $E_H=U_H/d$; setzt man dies oben ein und löst nach
$U_H$ auf,
\[
  \frac{U_H}{d}=vB \quad\Longrightarrow\quad U_H=Bvd,
\]
erhält man die \textbf{\emph{Hall-Spannung}}. Dabei ist $d$ der Abstand
der Messelektroden bzw.\ der Durchmesser des Leiters oder Gefässes. Sind
$U_H$, $B$ und $d$ bekannt, lässt sich daraus umgekehrt die Geschwindigkeit
\[
  v=\frac{U_H}{Bd}
\]
der Ladungsträger berechnen.

\textbf{Medizinische Anwendung:} Blut enthält geladene Teilchen (Ionen im
Blutplasma), die sich mit dem Blutfluss bewegen. Legt man von aussen ein
Magnetfeld $B$ quer über ein Blutgefäss und misst mit zwei Elektroden auf
gegenüberliegenden Seiten die Hall-Spannung $U_H$, lässt sich bei bekanntem
Gefässdurchmesser $d$ die Fliessgeschwindigkeit des Blutes berechnen --
nicht-invasiv, ohne das Gefäss zu öffnen.

\Bild{7cm}{Bilder/blut_messung.png}

\textbf{Aufgabe zum Vertiefen:}
\href{https://www.leifiphysik.de/elektrizitaetslehre/bewegte-ladungen-feldern/aufgabe/medizinische-anwendung-des-hall-effekts-abitur-2013-ph11-a2-1}{Medizinische
Anwendung des Hall-Effekts (Abitur BY 2013 Ph11 A2-1)}.
\end{theoriebox}

\begin{questions}
\setcounter{question}{4}
\begin{taskbox}[Kontrollfrage: Hall-Effekt (Gruppe E)]
\question[5]
\begin{parts}
  \part[3] Wie lässt sich mit dem Hall-Effekt die Fliessgeschwindigkeit
  von Blut messen, ohne das Gefäss zu öffnen? Welche zwei Grössen müssen
  dafür zusätzlich zur gemessenen Hall-Spannung $U_H$ bekannt sein?
  \begin{answerbox}
  \vspace{2.2cm}
  \end{answerbox}
  \begin{solution}
  Legt man ein bekanntes Magnetfeld $B$ quer über das Gefäss (Durchmesser
  $d$) und misst mit zwei Elektroden die Hall-Spannung $U_H$, lässt sich
  die Fliessgeschwindigkeit über $v=U_H/(Bd)$ berechnen, ohne das Gefäss
  zu öffnen. Zusätzlich zur gemessenen Hall-Spannung $U_H$ müssen das
  Magnetfeld $B$ und der Gefässdurchmesser $d$ bekannt sein.
  \end{solution}

  \part[2] Über einem stromdurchflossenen Metallleiter der Breite
  $d=\SI{2.0}{\centi\meter}$ liegt ein Magnetfeld $B=\SI{0.80}{\tesla}$.
  Die Ladungsträger bewegen sich mit einer typischen
  Driftgeschwindigkeit $v=\SI{1.5e-3}{\meter\per\second}$. Berechnen Sie
  die entstehende Hall-Spannung $U_H$.
  \begin{answerbox}
  \vspace{1.8cm}
  \end{answerbox}
  \begin{solution}
  $U_H=Bvd=\SI{0.80}{\tesla}\cdot\SI{1.5e-3}{\meter\per\second}\cdot\SI{2.0e-2}{\meter}
  =\SI{2.4e-5}{\volt}=\SI{24}{\micro\volt}$ -- deutlich kleiner als die
  Hall-Spannung im Blutgefäss oben, weil die Driftgeschwindigkeit in einem
  Metall viel kleiner ist als eine typische Blutflussgeschwindigkeit.
  \end{solution}
\end{parts}
\end{taskbox}
\end{questions}

% =============================================================================
% Vertiefungsphase: Stammgruppe wieder komplett
% =============================================================================
\textbf{Vertiefungsphase: Ihre Stammgruppe ist wieder komplett}

Nachdem in Ihrer Stammgruppe alle fünf Themen einmal erklärt wurden, lösen
alle fünf Mitglieder gemeinsam die folgenden zwei Aufgaben -- jede Aufgabe
braucht Wissen aus mindestens zwei verschiedenen Expertenthemen.

\needspace{14cm}
\begin{questions}
\setcounter{question}{5}
\begin{groupbox}[Stammgruppenaufgabe: Massenspektrograph und Hall-Effekt]
\question[7] Beide Themen beruhen auf demselben Kräftegleichgewicht bei
gekreuzten Feldern ($qE=qvB$ bzw.\ $qE_H=qvB$), nutzen es aber in
entgegengesetzter Richtung: Beim Wien-Filter (Stufe 1 des
Massenspektrographen) wird $E$ so eingestellt, dass genau eine
\emph{vorgegebene} Geschwindigkeit $v_0$ selektiert wird; beim Hall-Effekt
wird umgekehrt aus einer \emph{gemessenen} Spannung $U_H$ eine zunächst
\emph{unbekannte} Geschwindigkeit $v$ berechnet.
\begin{parts}
  \part[3] Ein Ionenstrahl durchläuft einen Wien-Filter mit $E=\SI{1.2e5}{\volt\per\meter}$
  und $B_1=\SI{0.40}{\tesla}$ und anschliessend ein zweites Magnetfeld
  $B_2=\SI{0.60}{\tesla}$, in dem ein Halbkreis mit Radius
  $r=\SI{0.125}{\meter}$ gemessen wird. Bestimmen Sie mit
  $v_0=E/B_1$ und $r=mv_0/qB_2 \Rightarrow m=qB_2r/v_0$ (Ladung
  $q=\SI{1.6e-19}{\coulomb}$) die Masse $m$ des Ions.
  \begin{answerbox}
  \vspace{3cm}
  \end{answerbox}
  \begin{solution}
  $v_0=E/B_1=\SI{1.2e5}{\volt\per\meter}/\SI{0.40}{\tesla}=\SI{3.0e5}{\meter\per\second}$.\\
  $m=\dfrac{qB_2r}{v_0}=\dfrac{\SI{1.6e-19}{\coulomb}\cdot\SI{0.60}{\tesla}\cdot\SI{0.125}{\meter}}{\SI{3.0e5}{\meter\per\second}}
  =\SI{4.0e-26}{\kilogram}\approx\SI{24.1}{u}$ -- das entspricht etwa der
  Masse eines Magnesium-Ions ($^{24}$Mg).
  \end{solution}

  \part[2] Über ein Blutgefäss mit Durchmesser $d=\SI{4.0}{\milli\meter}$
  wird ein Magnetfeld $B=\SI{0.50}{\tesla}$ gelegt und eine Hall-Spannung
  $U_H=\SI{1.0}{\milli\volt}$ gemessen. Bestimmen Sie mit
  $U_H=Bvd \Rightarrow v=U_H/(Bd)$ die Fliessgeschwindigkeit $v$ des
  Blutes.
  \begin{answerbox}
  \vspace{2.2cm}
  \end{answerbox}
  \begin{solution}
  $v=\dfrac{U_H}{Bd}=\dfrac{\SI{1.0e-3}{\volt}}{\SI{0.50}{\tesla}\cdot\SI{4.0e-3}{\meter}}
  =\SI{0.50}{\meter\per\second}$ -- eine realistische
  Blutflussgeschwindigkeit in einem grösseren Gefäss.
  \end{solution}

  \part[2] Diskutieren Sie zu zweit (Massenspektrograph- und
  Hall-Effekt-Experte/in): Was ist an den beiden Rechnungen strukturell
  gleich, und was genau ist vertauscht -- welche Grösse ist jeweils
  gegeben, welche gesucht?
  \begin{answerbox}
  \vspace{2.2cm}
  \end{answerbox}
  \begin{solution}
  Beide Rechnungen beruhen auf demselben Kräftegleichgewicht $qE=qvB$ bei
  gekreuzten Feldern. Beim Wien-Filter ist die Geschwindigkeit $v_0$ über
  $E$ und $B_1$ \emph{vorgegeben} (gesucht ist dort in Stufe 2 die Masse
  $m$); beim Hall-Effekt ist umgekehrt $B$ bekannt und die Geschwindigkeit
  $v$ \emph{gesucht}, berechnet aus der gemessenen Spannung $U_H$. Es ist
  also vertauscht, ob die Geschwindigkeit die gegebene oder die gesuchte
  Grösse ist.
  \end{solution}
\end{parts}
\end{groupbox}
\end{questions}

\needspace{12cm}
\begin{questions}
\setcounter{question}{6}
\begin{groupbox}[Stammgruppenaufgabe: Synthesetabelle]
\question[4] Füllen Sie gemeinsam die Tabelle aus: eine Zeile pro Thema,
mit den beteiligten Feldern, der Rolle der Lorentzkraft und dem Zweck der
Anwendung. Diskutieren Sie danach zu fünft: Was ist allen fünf Anwendungen
gemeinsam?
\begin{center}
\small
\begin{tabularx}{\textwidth}{@{}P{3.2cm}P{2.9cm}P{3.7cm}Y@{}}
\toprule
\textbf{Thema} & \textbf{Beteiligte Felder} & \textbf{Rolle der Lorentzkraft} & \textbf{Zweck} \\
\midrule
Massenspektrograph & & & \\[1.4cm]
Kreisbeschleuniger & & & \\[1.4cm]
MHD-Antrieb & & & \\[1.4cm]
Polarlichter & & & \\[1.4cm]
Hall-Effekt & & & \\[1.4cm]
\bottomrule
\end{tabularx}
\end{center}
\begin{solution}
\begin{center}
\small
\begin{tabularx}{\textwidth}{@{}P{3.2cm}P{2.9cm}P{3.7cm}Y@{}}
\toprule
\textbf{Thema} & \textbf{Beteiligte Felder} & \textbf{Rolle der Lorentzkraft} & \textbf{Zweck} \\
\midrule
Massenspektrograph & $E$ und $B$ gekreuzt (Stufe 1), nur $B$ (Stufe 2) &
Kräftegleichgewicht (Stufe 1), Zentripetalkraft (Stufe 2) & Massen (Isotope)
trennen \\[0.3cm]
Kreisbeschleuniger & nur $B$ & Zentripetalkraft & Teilchen beschleunigen \\[0.3cm]
MHD-Antrieb & nur $B$ (plus Strom $I$) & direkte Antriebskraft & Wasser
antreiben \\[0.3cm]
Polarlichter & nur $B$ & Zentripetalkraft (Spiralbahn) & Teilchen in die
Atmosphäre lenken \\[0.3cm]
Hall-Effekt & $E_H$ und $B$ gekreuzt & Kräftegleichgewicht &
Geschwindigkeit messen \\
\bottomrule
\end{tabularx}
\end{center}

Gemeinsam ist allen fünf Anwendungen, dass dieselbe Lorentzkraft $F=qvB$
(bzw.\ $F=BIL$ über den Strom) je nach Situation entweder als
Zentripetalkraft eine Kreis- oder Spiralbahn erzeugt, als Teil eines
Kräftegleichgewichts bei gekreuzten Feldern eine Spannung bzw.\ eine feste
Geschwindigkeit liefert, oder als direkte Antriebskraft wirkt.
\end{solution}
\end{groupbox}
\end{questions}

\end{document}
